
Debiasing methods such as Just Train Twice (JTT) succeed by reweighting examples that are learned later in training, implicitly relying on temporal ordering between spurious and core features. While JTT validates this operationally, the mechanistic foundation---when features converge at the gradient level---has not been measured. Without direct gradient-level verification, the distinction between temporal ordering and other factors (example hardness, loss landscape geometry) cannot be established, and principled interventions exploiting learning dynamics cannot be designed.

Standard training (ERM) on spurious correlation benchmarks like Waterbirds achieves 97\% accuracy on majority groups but collapses to 40\% on minority groups where spurious cues mislead the model. Existing debiasing methods---GroupDRO, Invariant Risk Minimization, and JTT---improve worst-group accuracy empirically but lack mechanistic grounding in optimization dynamics. The temporal hypothesis---that spurious features converge earlier than core features---is stated in JTT and Learning from Failure papers but has never been verified via gradient-level measurement.

The gap is both conceptual and methodological. Conceptually, direct evidence that gradient descent's implicit bias toward simpler features manifests as temporal separation between spurious and core feature convergence is lacking. Methodologically, measuring when features (not examples) converge requires isolating feature types---spurious-only vs core-only training variants---then tracking per-epoch gradient norms until stabilization. No prior work has performed this measurement.

This work hypothesizes that spurious features converge earlier measurably via per-epoch gradient norm tracking on ablated feature types. If spurious features (color, background) converge at epoch $E_s$ and core features (shape, semantics) converge at $E_c$ with $E_c - E_s \geq 2$ epochs, gradient descent's implicit bias toward simpler features creates an exploitable temporal gap arising from the feature complexity hierarchy: simpler low-level features (color, texture) stabilize in early convolutional layers before higher-level semantic features (shape, morphology) requiring deeper hierarchical processing.

The following contributions are made:

\textbf{1. Spurious features converge 4 epochs earlier than core features} ($\Delta=4$, exceeding threshold by $2\times$ in single-dataset proof-of-concept), providing the first direct gradient-level validation of the temporal ordering hypothesis underlying JTT/LfF. This mechanistic evidence explains why reweighting late-learned examples succeeds: it implicitly corrects for temporal misalignment driven by gradient descent's bias toward simpler features.

\textbf{2. Early convolutional layers drive temporal gaps} via significantly higher spurious correlation than late layers (mean $\rho_j=0.004$ vs $0.001, p=0.0028, t=2.78$, Cohen's $d=0.25$), confirming architectural feature hierarchy as mechanism. Layer-wise gradient analysis validates that temporal separation arises from hierarchical processing, not dataset artifacts.

\textbf{3. Cross-dataset transfer fails catastrophically} (39\% vs 86\% target worst-group accuracy), revealing that neuron correlations cannot be reused across spurious feature types---a fundamental constraint for gradient-aware methods. This negative result identifies the theoretical boundary: neuron-level interventions require per-dataset calibration, limiting scalability compared to example-level reweighting.

\textbf{4. Feature-level forgetting rate} provides independent stability metric (51\% lower for spurious features), confirming that spurious features converge to more stable predictions consistent with simpler decision boundaries.

Temporal ordering is measurable, mechanistically grounded in architectural feature hierarchy, and explains why reweighting methods (JTT, LfF) succeed: they implicitly correct for temporal misalignment by upweighting late-learned examples containing core features. However, neuron-level gradient modulation faces dataset-specificity constraints---neuron correlation values cannot be reused across spurious types, limiting the scalability of fine-grained interventions.

\textbf{Scope and limitations.} Results are validated on CMNIST color-based spurious correlation (single-dataset proof-of-concept with seed 0; full 10-seed statistical validation in progress). Generalization to background (Waterbirds), attribute (CelebA), and context (NICO++) spurious types is future work. Multi-dataset validation requires dataset-specific ablation strategies but uses identical gradient convergence measurement protocol.
